
\section{Conclusion}
\label{sec:conclusion}

The standard formulation of tool routing asks the model to predict tools.
Our results suggest this is the wrong prediction target.
When the model predicts entity types instead and graph search handles composition, routing becomes simpler, more reliable, and more scalable---consistently outperforming direct tool selection across four language models, five benchmark domains, and a production MCP server with over one thousand tools.
The decomposition eliminates hallucinated tool invocations by construction, and the recall decomposition reveals that end-to-end performance separates cleanly into entity type prediction quality and graph reachability---two independent and actionable levers for improvement.

The decomposition is the primary contributor to performance gains, not the choice of predictor.
TCR improves every model--domain combination we tested.
A compact domain-specific classifier ($\sim$475K trainable parameters) matches Gemini 2.5 Pro with few-shot prompting on two production-scale APIs spanning 536 to 1,060 tools, confirming that the semantic component can be served by lightweight, domain-specific models rather than general-purpose frontier models.

TCR assumes that the tool ecosystem exposes an underlying typed transformation graph.
Where this assumption holds---typed API ecosystems, infrastructure automation, enterprise services---the decomposition provides structural guarantees that no amount of model scaling can replicate.
Where it does not, the method does not apply.
Identifying the boundary between these regimes, and extending the formulation to domains with weaker type structure, remains an open direction.

More broadly, improving tool routing may require changing the prediction target---from tools to the domain abstractions that organize them---rather than building increasingly capable tool-selecting language models.
